% Introducción íntegra de EGA I, pp.5--9; no es el tomo completo. Fragmento editable para el proyecto del tomo.

% Unidad R565: introducción, p. 5, dedicatoria y requisitos a)--d).
% Continuación: intro-fr.tex, párrafo «Ce n'est pas le lieu...».
\oldpage[I]{5}
\section*{Introducción}

\begin{center}
\emph{A Oscar Zariski y André Weil.}
\end{center}


Este trabajo, y los numerosos trabajos que han de seguirlo, están destinados a formar un tratado sobre los fundamentos de la geometría algebraica.
En principio no presuponen ningún conocimiento particular de esta disciplina; incluso se ha comprobado que tal conocimiento, pese a sus evidentes ventajas, podía a veces resultar perjudicial para quien desea familiarizarse con el punto de vista y las técnicas expuestos aquí, por el hábito de adoptar de manera demasiado exclusiva el punto de vista birracional que dicho conocimiento lleva consigo.
En cambio, supondremos que el lector conoce bien los temas siguientes:

\begin{enumerate}
  \item[\textit{a)}] El \emph{álgebra conmutativa}, tal como se expone, por ejemplo, en los volúmenes de los \emph{Éléments} de N.~Bourbaki que están en preparación (y, hasta que aparezcan esos volúmenes, en Samuel-Zariski~\cite{I-13} y Samuel~\cite{I-11},~\cite{I-12}).
  \item[\textit{b)}] El \emph{álgebra homológica}, para la cual remitimos a Cartan-Eilenberg~\cite{I-2} (citado como~(M)) y Godement~\cite{I-4} (citado como~(G)), así como al reciente artículo de A.~Grothendieck~\cite{I-6} (citado como~(T)).
  \item[\textit{c)}] La \emph{teoría de haces}, cuyas referencias principales serán~(G) y~(T); esta teoría proporciona el lenguaje indispensable para interpretar en términos «geométricos» las nociones esenciales del álgebra conmutativa y para «globalizarlas».
  \item[\textit{d)}] Por último, será útil que el lector esté familiarizado en cierta medida con el \emph{lenguaje funtorial}, que se empleará constantemente en este tratado y para el cual podrá consultar~(M),~(G) y, sobre todo,~(T); los principios de este lenguaje y los principales resultados de la teoría general de los funtores se expondrán con mayor detalle en una obra que los autores de este tratado están preparando.
\end{enumerate}

\bigskip
\begin{center}*\quad*\quad*\end{center}
\bigskip

% Unidad R566: introducción, líneas francesas 24--35, pp. 5--6.
% Traducción: OpenAI Codex, GPT-6.1 Sol, esfuerzo Ultra.
No corresponde en esta Introducción ofrecer una descripción más o menos sumaria del punto de vista de los «esquemas» en geometría algebraica, ni la larga lista de razones que hicieron necesaria su adopción, en particular la aceptación sistemática de elementos nilpotentes en los anillos locales de las «variedades» que consideramos (lo que necesariamente relega a un segundo plano la noción de aplicación racional, en favor de la de aplicación regular o «morfismo»).
Este tratado se propone precisamente desarrollar de manera sistemática el lenguaje de los «esquemas» y demostrará, así lo esperamos, su necesidad.
Aunque sería fácil hacerlo, tampoco intentaremos ofrecer aquí

\oldpage[I]{6}
una introducción «intuitiva» a las nociones desarrolladas en el capítulo primero.
El lector que desee tener una visión preliminar de las materias de este tratado podrá consultar la conferencia de A.~Grothendieck en el Congreso Internacional de Matemáticos celebrado en Edimburgo en 1958~\cite{I-7}, así como la exposición~\cite{I-8} del mismo autor.
El trabajo~\cite{I-14} (citado como~(FAC)) de J.-P.~Serre también puede considerarse una exposición intermedia entre el punto de vista clásico y el de los esquemas en geometría algebraica; por ello, su lectura puede constituir una excelente preparación para la de nuestros \emph{Éléments}.

\bigskip
\begin{center}*\quad*\quad*\end{center}
\bigskip

% Unidad R567: introducción, líneas francesas 37--68, pp. 6--7.
% Traducción: OpenAI Codex, GPT-6.1 Sol, esfuerzo Ultra.
A título informativo, damos a continuación el plan general previsto para este tratado, sujeto, por lo demás, a modificaciones posteriores, sobre todo en lo que respecta a los últimos capítulos:

\medskip
\begin{tabular}{@{}p{0.22\linewidth}p{0.74\linewidth}@{}}
  Capítulo primero. & --- El lenguaje de los esquemas.\\
  --- II. & --- Estudio global elemental de algunas clases de morfismos.\\
  --- III. & --- Cohomología de los haces algebraicos coherentes. Aplicaciones.\\
  --- IV. & --- Estudio local de los morfismos.\\
  --- V. & --- Procedimientos elementales de construcción de esquemas.\\
  --- VI. & --- Técnica de descenso. Método general de construcción de esquemas.\\
  --- VII. & --- Esquemas de grupos, espacios fibrados principales.\\
  --- VIII. & --- Estudio diferencial de los espacios fibrados.\\
  --- IX. & --- El grupo fundamental.\\
  --- X. & --- Residuos y dualidad.\\
  --- XI. & --- Teorías de intersección, clases de Chern, teorema de Riemann-Roch.\\
  --- XII. & --- Esquemas abelianos y esquemas de Picard.\\
  --- XIII. & --- Cohomología de Weil.\\
\end{tabular}
\medskip

En principio, todos los capítulos se consideran abiertos a ampliaciones, y siempre podrán añadírseles párrafos suplementarios más adelante; estos párrafos aparecerán en fascículos separados, para reducir los inconvenientes del modo de publicación adoptado.
Cuando uno de estos párrafos esté previsto o en preparación al publicarse un capítulo, se mencionará en el sumario de ese capítulo, aunque determinadas prioridades de urgencia hagan que su publicación efectiva tenga lugar mucho más tarde.
Para comodidad del lector, ofrecemos en un «Capítulo~0» diversos complementos de álgebra conmutativa, álgebra homológica y teoría de haces, utilizados en los capítulos de este tratado, que son más o menos conocidos, pero para los cuales no ha sido posible dar referencias cómodas de consultar.
Se recomienda al lector acudir al capítulo~0 solo mientras lee el tratado propiamente dicho, y en la medida en que los resultados a los que remitimos no le sean

\oldpage[I]{7}
suficientemente familiares.
Pensamos, además, que así la lectura de este tratado podrá ser para el principiante un buen método de familiarizarse con el álgebra conmutativa y el álgebra homológica, cuyo estudio, cuando no se acompaña de aplicaciones tangibles, resulta tedioso, e incluso deprimente, para un número considerable de personas.

\bigskip
\begin{center}*\quad*\quad*\end{center}
\bigskip

% Unidad R568: introducción, líneas francesas 70--93, pp. 7--8.
% Traducción: OpenAI Codex, GPT-6.1 Sol, esfuerzo Ultra.
Queda fuera de nuestra competencia ofrecer en esta introducción una reseña histórica, siquiera sumaria, de las nociones y los resultados expuestos.
El texto solo contendrá las referencias que se consideren especialmente útiles para su comprensión, y solo indicaremos el origen de los resultados más importantes.
Al menos en su forma, los temas tratados en nuestra obra son bastante nuevos; esto explica la escasez de referencias a los padres de la geometría algebraica del siglo~\textsc{xix} y de comienzos del siglo~\textsc{xx}, cuyos trabajos solo conocemos de oídas.
Conviene, sin embargo, decir aquí unas palabras sobre las obras que han influido más directamente en los autores y han contribuido al desarrollo del punto de vista de los esquemas.
Hay que citar, en primer lugar, el trabajo fundamental (FAC) de J.-P.~Serre, que ha servido de introducción a la geometría algebraica a más de un joven interesado en ella (entre ellos, uno de los autores de este tratado), desalentado por la aridez de los clásicos \emph{Foundations} de A.~Weil~\cite{I-18}.
Allí se demuestra por primera vez que la «topología de Zariski» de una variedad algebraica «abstracta» es plenamente adecuada para aplicarle ciertas técnicas de la topología algebraica y dar lugar, en particular, a una teoría cohomológica.
Además, la definición de variedad algebraica que allí se da es la que se presta de manera más natural a la extensión de esta noción que desarrollamos aquí\footnote{Como nos ha señalado J.-P.~Serre, conviene recordar que la idea de definir la estructura de variedad mediante un haz de anillos se debe a H.~Cartan, quien tomó esta idea como punto de partida de su teoría de los espacios analíticos.
Por supuesto, al igual que en geometría algebraica, sería importante admitir, en «geometría analítica», los elementos nilpotentes en los anillos locales de los espacios analíticos.
H.~Grauert ha abordado recientemente esta extensión de la definición de H.~Cartan y J.-P.~Serre~\cite{I-5}, y cabe esperar que pronto aparezca una exposición sistemática de geometría analítica en este marco general.
Es evidente, por otra parte, que las nociones y técnicas desarrolladas en este tratado siguen teniendo sentido en geometría analítica, aunque en esta última teoría hay que prever dificultades técnicas más considerables.
Cabe prever que la geometría algebraica, gracias a la sencillez de sus métodos, podrá servir como una especie de modelo formal para futuros desarrollos de la teoría de los espacios analíticos.}.
El propio Serre había observado, además, que la teoría cohomológica de las variedades algebraicas afines podía trasladarse sin dificultad sustituyendo las álgebras afines sobre un cuerpo por anillos conmutativos cualesquiera.
Por tanto, los capítulos~I y~II de este tratado y los dos primeros párrafos del capítulo~III pueden considerarse, en lo esencial, transposiciones fáciles, a este marco ampliado, de los principales resultados de (FAC) y de un artículo posterior del mismo autor~\cite{I-15}.
También hemos obtenido gran provecho del \emph{Séminaire de Géométrie algébrique} de C.~Chevalley~\cite{I-1}; en particular, el uso sistemático de los «conjuntos construibles» introducidos por él ha resultado muy útil en la teoría de esquemas (cf. cap.~IV).
También hemos tomado de él el estudio de los morfismos desde el punto

\oldpage[I]{8}
de vista de la dimensión (cap.~IV), que se traslada al marco de los esquemas sin cambios apreciables.
Conviene señalar, además, que la noción de «esquemas de anillos locales», introducida por Chevalley, se presta de manera natural a una extensión de la geometría algebraica (que, sin embargo, no tiene toda la flexibilidad y generalidad que pretendemos darle aquí); para la relación entre esta noción y nuestra teoría, véase el capítulo primero, \S~8.
M.~Nagata ha desarrollado una extensión de este tipo en una serie de memorias~\cite{I-9} que contienen numerosos resultados especiales sobre geometría algebraica sobre anillos de Dedekind\footnote{Entre los trabajos que se aproximan a nuestro punto de vista en geometría algebraica, señalemos el importante trabajo de E.~Kähler~\cite{I-22} y una nota reciente de Chow e Igusa~\cite{I-3}, que retoman, en el marco de la teoría de Nagata-Chevalley, ciertos resultados de (FAC) y dan también una fórmula de Künneth.}.

\bigskip
\begin{center}*\quad*\quad*\end{center}
\bigskip

% Unidad R569: introducción, líneas francesas 95--120, pp. 8--9.
% Traducción: OpenAI Codex, GPT-6.1 Sol, esfuerzo Ultra.
Por último, es evidente que un libro sobre geometría algebraica, y sobre todo un libro sobre sus fundamentos, recibe necesariamente la influencia, aunque sea por intermediarios, de matemáticos como O.~Zariski y A.~Weil.
En particular, la \emph{Théorie des fonctions holomorphes} de Zariski~\cite{I-20}, convenientemente flexibilizada gracias a los métodos cohomológicos y completada con un teorema de existencia (cap.~III, \S\S~4 y~5), es, junto con la técnica de descenso expuesta en el cap.~VI, una de las principales herramientas empleadas en este tratado y nos parece una de las más potentes de que se dispone en geometría algebraica.

La técnica general en la que se inscribe puede esbozarse del siguiente modo (se ofrecerá un ejemplo típico en el cap.~IX, al estudiar el grupo fundamental).
Se tiene un morfismo propio (cap.~II) $f:X\to Y$ de una variedad algebraica a otra (más generalmente, de un esquema a otro) que se quiere estudiar en una vecindad de un punto $y\in Y$, con el fin de resolver un problema P relativo a una vecindad de $y$.
Se procede por etapas sucesivas:
\begin{enumerate}
  \item[1\textsuperscript{o}] Se puede suponer que $Y$ es afín, de modo que $X$ pasa a ser un esquema definido sobre el anillo afín $A$ de $Y$, e incluso se puede sustituir $A$ por el anillo local de $y$.
  Esta reducción siempre es fácil en la práctica (cap.~V) y nos lleva al caso en que $A$ es un anillo \emph{local}.
  \item[2\textsuperscript{o}] Se estudia el problema considerado cuando $A$ es un anillo local \emph{artiniano}.
  Para que siga teniendo efectivamente sentido cuando no se supone que $A$ sea íntegro, a veces es preciso reformular el problema P, y resulta que a menudo se obtiene así una mejor comprensión del problema, de naturaleza «infinitesimal» en esta etapa.
  \item[3\textsuperscript{o}] La teoría de esquemas formales (cap.~III, \S\S~3, 4 y~5) permite pasar del caso de un anillo artiniano al de un \emph{anillo local completo}.
  \item[4\textsuperscript{o}] Por último, si $A$ es un anillo local cualquiera, la consideración de «secciones multiformes» sobre esquemas adecuados sobre $X$, que aproximan una sección «formal» dada (cap.~IV), permitirá a menudo pasar de un resultado conocido para el esquema

  \oldpage[I]{9}
  deducido de $X$ por extensión de escalares a la completación de $A$, a un resultado análogo para una extensión finita suficientemente sencilla (por ejemplo, no ramificada) de $A$.
\end{enumerate}

Este esbozo muestra la importancia del estudio sistemático de los esquemas definidos sobre un anillo artiniano $A$.
El punto de vista de Serre en su formulación de la teoría local de cuerpos de clases, y algunos trabajos recientes de Greenberg, parecen sugerir que ese estudio podría emprenderse asociando funtorialmente a un esquema $X$ de este tipo un esquema $X'$ sobre el cuerpo residual $k$ de $A$ (supuesto perfecto), de dimensión igual, en los casos favorables, a $n\dim X$, donde $n$ es la longitud de $A$.

En cuanto a la influencia de A.~Weil, baste decir que la necesidad de desarrollar las herramientas para formular con toda la generalidad deseada la definición de la «cohomología de Weil» y para abordar la demostración\footnote{Para evitar cualquier malentendido, precisemos que esta tarea apenas acaba de emprenderse en el momento en que se escribe esta introducción y, por tanto, todavía no ha llevado a la demostración de las conjeturas de Weil.} de todas las propiedades formales necesarias para establecer sus célebres conjeturas en geometría diofántica~\cite{I-19} ha sido una de las principales motivaciones para redactar este tratado, al igual que el deseo de encontrar el marco natural de las nociones y métodos habituales de la geometría algebraica y de dar a los autores la oportunidad de comprender esas nociones y técnicas.

\bigskip
\begin{center}*\quad*\quad*\end{center}
\bigskip

% Unidad R570: introducción, líneas francesas 122--132, p. 9.
% Traducción: OpenAI Codex, GPT-6.1 Sol, esfuerzo Ultra.
Para terminar, creemos útil advertir a los lectores que, al igual que los propios autores, seguramente encontrarán alguna dificultad antes de acostumbrarse al lenguaje de los esquemas y de convencerse de que las construcciones habituales sugeridas por la intuición geométrica pueden trasladarse a este lenguaje, esencialmente de una única manera razonable.
Como en muchas partes de la matemática moderna, la intuición inicial parece alejarse cada vez más del lenguaje adecuado para expresarla con toda la precisión y generalidad deseadas.
En este caso, la dificultad psicológica se debe a la necesidad de trasladar a los objetos de una categoría ya bastante distinta de la categoría de conjuntos (a saber, la categoría de preesquemas o la de preesquemas sobre un preesquema dado) nociones familiares para los conjuntos: productos cartesianos, leyes de grupo, de anillo, de módulo, fibrados, fibrados principales homogéneos, etc.
En el futuro, seguramente será difícil para el matemático eludir este nuevo esfuerzo de abstracción, quizá bastante pequeño, después de todo, en comparación con el que hicieron nuestros predecesores al familiarizarse con la teoría de conjuntos.

\bigskip
\begin{center}*\quad*\quad*\end{center}
\bigskip

Las referencias se darán según el sistema decimal; por ejemplo, en III, 4.9.3, III indica el capítulo, 4 el párrafo y 9 la sección del párrafo.
Dentro \emph{de un mismo capítulo}, se omitirá la indicación del capítulo.

